% Retained supplementary source fragment. This is not a replacement paper.
% It uses the parent manuscript's theorem environments and author-date references.
\section{Supplementary proof for signed Gaussian loading and its competitors}
\label{sec:signed-gaussian-loading-supplement}

I retain the established classical probability model throughout this
supplement. Quantum states encode computational proposals; they do not
replace the financial risk-neutral measure. I give the continuous
Gaussian identities first, then a finite-register identity which avoids
assuming that a translated continuous Gaussian is an exact modular
translation on a finite grid. I attribute the LCU, rejection and
amplification mechanisms to their established sources
\citep{rosenkranz2025,ozols2013,brassard}. The financial choice of the
remainder proposal and the comparisons below are derivations for this
project. I make no theoretical-priority or complete-cost claim.

\subsection{Probability measures, centring and price notation}

Let $p\geq1$, let $w_i>0$ with $\sum_{i=1}^n w_i=1$, and let
$d_i\in\mathbb R$ and $v_i\in\mathbb R^p$ satisfy
\begin{equation}
 \sum_iw_i d_i=0,\qquad \sum_iw_i v_i=0.
 \label{eq:sg-centering}
\end{equation}
Define
\[
 T(y)=\sum_iw_i e^{d_i+v_i\cdot y},\qquad U(y)=T(y)-1.
\]
Weighted Jensen gives $T(y)\geq1$ for every $y$. These are the centred
log-price contrasts obtained by dividing a positive arithmetic basket
by its weighted geometric basket. They are not a new asset-price model.
The continuous base law is the standard Gaussian $\gamma_p$, with density
\[
 \varphi_p(y)=(2\pi)^{-p/2}e^{-\|y\|^2/2}.
\]
I also use a fixed finite base law $\mu$, with masses $p_y>0$ on a
finite grid $\mathcal G$. Every finite identity below concerns that
same law, rather than substituting a continuous normaliser.

Fix an integer $s\geq1$ and put $r=2s$. For either base law $\rho$,
write
\[
 M_-^\rho=\mathbb E_\rho U^r,\qquad
 M_+^\rho=\mathbb E_\rho(T+1)^r.
\]
When I write $p_\rho$, it denotes $\varphi_p$ for $\gamma_p$ and
the mass $p_y$ for $\mu$. The corresponding proposal measures are
$q_-^\rho(dy)=U(y)^r\rho(dy)/M_-^\rho$ and
$q_+^\rho(dy)=(T(y)+1)^r\rho(dy)/M_+^\rho$.
Gaussian exponential moments imply finiteness of every integer
moment of $T$ under $\gamma_p$. All moments are finite under $\mu$.
If $M_-^\rho=0$, nonnegativity implies $U=0$ almost surely. Every
Taylor remainder containing the factor $U^r$ is then zero almost surely,
and no importance state is needed. I assume $M_-^\rho>0$ wherever
I normalise such a state.

Let $g_y(t)$ be the conditioned monetary payoff at ratio $t>0$.
In the capped GBM setting it is
$P_{\lambda,H}(m_G(y)+\log t)$, with $K,H,\lambda>0$.
For $J_{r-1}(y)=\sum_{j=0}^{r-1}g_y^{(j)}(1)U(y)^j/j!$,
define $R_r(y)=g_y(T(y))-J_{r-1}(y)$. I assume the checked derivative
certificate supplies a strictly positive constant $D_r^+$ such that
\[
 |g_y^{(r)}(t)|\leq H D_r^+t^{-r},\qquad t>0.
\]
An exact rational outward upper bound on the joint norm is sufficient.
Put $b_s=H D_r^+/r!$ and $B_s^\rho=b_sM_-^\rho$.
Taylor's integral formula yields $|R_r(y)|\leq b_sU(y)^r$.
The lower-order derivative certificates and finite moments likewise
make each term of $J_{r-1}$ integrable.
I write $V_\rho=\mathbb E_\rho g_Y(T(Y))$ for the undiscounted
expectation and $\Pi_\rho=e^{-r_f\tau}V_\rho$ for its discounted quantity,
where $r_f$ is the risk-free rate and $\tau$ the financial maturity.
The continuous quantity $\Pi_{\gamma_p}$ is the classical model price.
The finite expectation $V_\mu$ and its discounted counterpart are
numerical targets. Transfer from $V_{\gamma_p}$ requires the separate
tail and grid certificate; a finite-grid change of the Gaussian law is
not automatically a new risk-neutral market model.

\subsection{Continuous Gaussian tilts and their nonorthogonal norm}

For a multi-index $\alpha\in\mathbb N_0^n$, write
$j=|\alpha|$, $\alpha!=\prod_i\alpha_i!$, $w^\alpha=\prod_iw_i^{\alpha_i}$,
$d_\alpha=\sum_i\alpha_i d_i$ and
$v_\alpha=\sum_i\alpha_i v_i$. On
$\mathcal I_s=\{\alpha:|\alpha|\leq s\}$, define
\[
 c_\alpha=(-1)^{s-j}\binom sj\frac{j!}{\alpha!}w^\alpha,
 \qquad
 \beta_\alpha=c_\alpha e^{d_\alpha+\|v_\alpha\|^2}.
\]
The zero multi-index is included, with $d_0=0$, $v_0=0$ and
$w^0=1$. For a real vector $a$, set
\[
 \psi_a(y)=(2\pi)^{-p/4}e^{-\|y-a\|^2/4}.
\]
It is a unit vector in $L^2(\mathbb R^p)$ and its squared modulus is
the density of $N(a,I_p)$.

\begin{proposition}[Continuous amplitude identity]
\label{prop:sg-continuous}
For the centred contrasts above,
\begin{equation}
 U(y)^s\psi_0(y)=\sum_{\alpha\in\mathcal I_s}
                         \beta_\alpha\psi_{2v_\alpha}(y).
 \label{eq:sg-continuous-amplitude}
\end{equation}
The target normalised wavefunction has squared modulus
$q_-(y)=U(y)^r\varphi_p(y)/M_-^{\gamma_p}$.
If $L_s^\gamma=\sum_{\alpha\in\mathcal I_s}|\beta_\alpha|$, then
$0<M_-^{\gamma_p}\leq(L_s^\gamma)^2$. There are
$\binom{n+s}{s}$ labelled components before exact merging.
\end{proposition}
\begin{proof}
The binomial and multinomial identities give
$U^s=\sum_{\alpha\in\mathcal I_s}c_\alpha
e^{d_\alpha+v_\alpha\cdot y}$.
Completing the square gives
\[
 e^{v\cdot y}\psi_0(y)
   =e^{\|v\|^2}\psi_{2v}(y),
\]
which proves equation~\eqref{eq:sg-continuous-amplitude}. Its squared
$L^2$ norm is $\mathbb E_{\gamma_p}U^r=M_-^{\gamma_p}$.
The triangle inequality and unit component norms give the upper bound.
The number of degree-$j$ multi-indices is $\binom{n+j-1}{j}$.
Summing over $0\leq j\leq s$ gives $\binom{n+s}{s}$, by the
identity $\binom{a}{j}+\binom{a}{j+1}=\binom{a+1}{j+1}$
applied successively.
\end{proof}

The translated components are generally not orthogonal. A second
completion of the square gives
\[
 \langle\psi_a,\psi_b\rangle
 =e^{-\|a-b\|^2/8}.
\]
Consequently the exact norm also has the Gram representation
\begin{equation}
 M_-^{\gamma_p}
  =\sum_{\alpha,\zeta\in\mathcal I_s}
       \beta_\alpha\beta_\zeta
       e^{-\|v_\alpha-v_\zeta\|^2/2}.
 \label{eq:sg-gaussian-gram}
\end{equation}
Indeed the product of the two Gaussian wavefunctions equals
$e^{-\|a-b\|^2/8}\varphi_p(y-(a+b)/2)$ and integrates to the stated
overlap. Substituting the coefficients into a Gram term gives
$c_\alpha c_\zeta
e^{d_\alpha+d_\zeta+\|v_\alpha+v_\zeta\|^2/2}$, precisely the
Gaussian exponential moment of that term in the product of the two
binomial expansions. This independently confirms the normaliser.
Discarding the cross terms or replacing the norm by
$\sum_\alpha|\beta_\alpha|^2$ is incorrect.

These are Hilbert-space identities. Translation
$f(y)\mapsto f(y-a)$ is unitary on $L^2$, but that fact does not give
a finite-qubit Gaussian loader, its controlled version or its cost.

\subsection{An exact finite-grid identity and an LCU interface}

For a finite law $\mu$, define the positive exponential sum
\[
 Z_\mu(v)=\sum_{y\in\mathcal G}p_y e^{v\cdot y}
\]
and the normalised vector
\[
 |\varphi_v^\mu\rangle
   =\sum_{y\in\mathcal G}
       \frac{\sqrt{p_y}e^{v\cdot y}}{\sqrt{Z_\mu(2v)}}|y\rangle.
\]
If $\mu$ is a product law on a product grid, then
$Z_\mu(v)=\prod_{a=1}^p\sum_{y_a}p_a(y_a)e^{v_a y_a}$.
This factorisation follows by distributing the finite sums; it does
not invoke continuous Gaussian integration.
Set
$\beta_\alpha^\mu=c_\alpha e^{d_\alpha}\sqrt{Z_\mu(2v_\alpha)}$
and $L_s^\mu=\sum_\alpha|\beta_\alpha^\mu|$. Exact substitution gives
\begin{equation}
 \sum_\alpha\beta_\alpha^\mu|\varphi_{v_\alpha}^\mu\rangle
    =\sum_y\sqrt{p_y}U(y)^s|y\rangle.
 \label{eq:sg-finite-amplitude}
\end{equation}
The normalisers have finite closed-sum expressions
\[
 M_\pm^\mu=
 \sum_{j=0}^{r}(\pm1)^{r-j}\binom rj
 \sum_{|\alpha|=j}\frac{j!}{\alpha!}w^\alpha
               e^{d_\alpha}Z_\mu(v_\alpha),
\]
where the upper sign denotes $M_+$ and the lower sign denotes $M_-$.
For the continuous law the corresponding exponential sum is
$Z_{\gamma_p}(v)=e^{\|v\|^2/2}$. These expressions can be evaluated
with outward arithmetic. The alternating $M_-$ expression can require
high precision; a strictly positive lower enclosure is necessary before
division or amplification.

\begin{proposition}[Finite LCU success]
\label{prop:sg-finite-lcu}
Suppose a stated unitary $P_\alpha$ prepares
$|\varphi_{v_\alpha}^\mu\rangle$ from zero, with private work returned
to zero and with the specified component phase. Let $P_\chi$ prepare
\[
 |\chi\rangle=\sum_{\alpha\in\mathcal I_s}
           \sqrt{|\beta_\alpha^\mu|/L_s^\mu}\,|\alpha\rangle.
\]
Padding labels have zero weight and a stated unitary extension.
Define
\[
 {\rm SELECT}=\sum_\alpha|\alpha\rangle\langle\alpha|
       \otimes {\rm sign}(\beta_\alpha^\mu)P_\alpha,\qquad
 \mathcal B=(P_\chi^\dagger\otimes I)
               {\rm SELECT}(P_\chi\otimes I).
\]
Projection onto the zero selector yields
$\sum_y\sqrt{p_y}U(y)^s|y\rangle/L_s^\mu$.
Its success probability is
\begin{equation}
 \eta_{\rm LCU}=M_-^\mu/(L_s^\mu)^2.
 \label{eq:sg-finite-success}
\end{equation}
Conditional on success, measuring the financial label yields
$q_-(y)=p_yU(y)^r/M_-^\mu$.
\end{proposition}
\begin{proof}
The selector preparation and its adjoint each contribute
$\sqrt{|\beta_\alpha^\mu|/L_s^\mu}$ to the zero-selector block.
The controlled phase supplies the coefficient sign, so the projected
vector is $(L_s^\mu)^{-1}\sum_\alpha
\beta_\alpha^\mu|\varphi_{v_\alpha}^\mu\rangle$.
Equation~\eqref{eq:sg-finite-amplitude} identifies that vector.
Its squared norm is $M_-^\mu/(L_s^\mu)^2$. Normalising and squaring
each financial amplitude gives the stated probability mass.
\end{proof}

The component phases in SELECT are relative phases. A component
preparation certified only up to an arbitrary global phase is insufficient
when that unitary is controlled inside this sum. Nor may I replace
$|\varphi_v^\mu\rangle$ by a periodic shift of a base Gaussian without
proving the finite-grid identity and its normalisation. The proposition
is an interface conditional on the component circuits, not a claim that
those circuits have been emitted for every financial input.

\subsection{Finite precision and cancellation}

\begin{proposition}[Coefficient and component precision]
\label{prop:sg-component-error}
Let $u=\sum_\alpha\beta_\alpha^\mu|\varphi_\alpha\rangle$, so
$\|u\|=\sqrt{M_-^\mu}$. Suppose normalised prepared components
$|\widetilde\varphi_\alpha\rangle$ satisfy
$\|\widetilde\varphi_\alpha-\varphi_\alpha\|\leq\delta_\alpha$,
with their relative phases fixed, and coefficients satisfy
$\sum_\alpha|\widetilde\beta_\alpha-\beta_\alpha^\mu|\leq\epsilon_\beta$.
Then for $\widetilde u=\sum_\alpha
\widetilde\beta_\alpha|\widetilde\varphi_\alpha\rangle$,
\[
 \|\widetilde u-u\|\leq
 e_u:=\epsilon_\beta+\sum_\alpha|\beta_\alpha^\mu|\delta_\alpha.
\]
If $e_u<\sqrt{M_-^\mu}$, then $\widetilde u\ne0$ and
\begin{equation}
 \left\|\frac{\widetilde u}{\|\widetilde u\|}
               -\frac{u}{\sqrt{M_-^\mu}}\right\|
       \leq \frac{2e_u}{\sqrt{M_-^\mu}}.
 \label{eq:sg-normalised-precision}
\end{equation}
For a norm-one observable, the error of its conditional mean is at most
$4e_u/\sqrt{M_-^\mu}$.
\end{proposition}
\begin{proof}
Split the vector difference as
$\sum_\alpha\beta_\alpha^\mu
(\widetilde\varphi_\alpha-\varphi_\alpha)
+\sum_\alpha(\widetilde\beta_\alpha-\beta_\alpha^\mu)
\widetilde\varphi_\alpha$ and apply the triangle inequality.
The reverse triangle inequality gives
$\|\widetilde u\|\geq\sqrt{M_-^\mu}-e_u>0$.
Insert $\widetilde u/\sqrt{M_-^\mu}$ between the normalised vectors.
The two resulting distances are
$|\|\widetilde u\|-\sqrt{M_-^\mu}|/\sqrt{M_-^\mu}$
and $\|\widetilde u-u\|/\sqrt{M_-^\mu}$; each is at most
$e_u/\sqrt{M_-^\mu}$.
For unit vectors $a,b$ and $\|F\|\leq1$, insertion of one bra or ket
at a time gives
$|\langle a,Fa\rangle-\langle b,Fb\rangle|\leq2\|a-b\|$.
This proves the last assertion.
\end{proof}

In particular, common component error $\delta$ gives
$e_u\leq\epsilon_\beta+L_s^\mu\delta$. Small target norm amplifies
preparation error as well as preparation effort.
This bound is about prepared columns. For repeated coherent use,
I also need a full operator-norm certificate for the emitted preparation.
If the operator errors of $P_\chi$ and SELECT are at most
$\zeta_\chi$ and $\zeta_{\rm sel}$, telescoping the three-unitary
product gives
$\|\widetilde{\mathcal B}-\mathcal B\|
\leq2\zeta_\chi+\zeta_{\rm sel}$.
For a block-diagonal SELECT, its error is the maximum error of its
controlled blocks, including their phases, because the selector
subspaces are orthogonal.

More generally, if
$\|\widetilde{\mathcal B}-\mathcal B\|\leq\zeta<\sqrt{\eta_{\rm LCU}}$,
projection is contractive and the same normalisation proof gives
a conditional state error at most $2\zeta/\sqrt{\eta_{\rm LCU}}$.
Multiplying the norm-one mean error by $B_s^\mu$ gives the monetary
bound
\[
 4B_s^\mu\zeta/\sqrt{\eta_{\rm LCU}}
       =4b_sL_s^\mu\sqrt{M_-^\mu}\,\zeta.
\]
This conditional bound does not, by itself, certify a complete
amplified pricing circuit.

\subsection{State-specific amplification without a false unit-success assumption}

Let $\mathsf G$ be the success projector of a preparation $\mathcal B$,
and put $|\Psi\rangle=\mathcal B|0\rangle$.
For $0<\eta=\|\mathsf G\Psi\|^2<1$, define orthonormal states
$|g\rangle=\mathsf G|\Psi\rangle/\sqrt\eta$ and
$|b\rangle=(I-\mathsf G)|\Psi\rangle/\sqrt{1-\eta}$, and let
$\theta=\arcsin\sqrt\eta$.
With
$S_0=I-2|0\rangle\langle0|$,
$S_{\mathsf G}=I-2\mathsf G$ and
$Q=-\mathcal B S_0\mathcal B^\dagger S_{\mathsf G}$,
direct multiplication on $\operatorname{span}\{g,b\}$ gives
\begin{equation}
 Q^a\mathcal B|0\rangle
   =\sin((2a+1)\theta)|g\rangle+
             \cos((2a+1)\theta)|b\rangle,\qquad a\in\mathbb N_0.
 \label{eq:sg-amplification}
\end{equation}
Indeed $-\mathcal B S_0\mathcal B^\dagger=2|\Psi\rangle
\langle\Psi|-I$, and the two-dimensional matrix of $Q$ in the ordered
basis $(g,b)$ is
$\bigl(\begin{smallmatrix}\cos2\theta&\sin2\theta\\
-\sin2\theta&\cos2\theta\end{smallmatrix}\bigr)$.
Multiplication and induction prove equation~\eqref{eq:sg-amplification}.
The endpoint $\eta=1$ needs no amplification.

The actual good probability after $a$ iterations is
$\kappa_a=\sin^2((2a+1)\theta)$, generally less than one.
For $\eta\geq1/2$, choose $a=0$. For $\eta<1/2$, choose the
nonnegative integer nearest $\pi/(4\theta)-1/2$.
Then $|(2a+1)\theta-\pi/2|\leq\theta$ and
$\kappa_a\geq\cos^2\theta>1/2$.
The number of preparation or adjoint occurrences is $2a+1$,
which is $O(\eta^{-1/2})$ since $\theta\geq\sqrt\eta$.
For this ordinary-amplification family, attaining any fixed
$\kappa_a\geq\kappa_0>0$ also requires
$2a+1\geq\sqrt{\kappa_0}/\theta$, using $|\sin x|\leq|x|$.
Thus the occurrence count is $\Theta(\eta^{-1/2})$ as $\eta\downarrow0$.
This is not a lower bound for every representation or preparation of
the target financial distribution.

In a numerical implementation, intervals for $\eta$ and $\theta$ must
certify one chosen integer $a$ and a strictly positive lower bound
for its $\kappa_a$. An uncertain rounding boundary cannot be treated
as an exact iteration decision. Generalised final phases can reduce
the failure probability, but their synthesis errors and gates must
also be retained. Oblivious amplification results for operators close
to unitary cannot be imported as a unit-success guarantee for this
non-unitary financial filter.

\begin{proposition}[A fully coherent monetary correction]
\label{prop:sg-coherent-correction}
Let $F_s(y)=R_r(y)/(b_sU(y)^r)$ when $U(y)>0$, and set it to any
bounded value where $U=0$. Then $|F_s|\leq1$ and
\begin{equation}
 V_\rho=\mathbb E_\rho J_{r-1}
                    +B_s^\rho\,\mathbb E_{q_-}F_s.
 \label{eq:sg-price-identity}
\end{equation}
After state-specific amplification, define the signed observable
$W=\mathsf G F_s$, taking value zero on failure. It has norm at most
one and mean $\kappa_a\mathbb E_{q_-}F_s$. A flag encoding
$(1+W)/2$ therefore has probability
$p_W=(1+\kappa_a\mathbb E_{q_-}F_s)/2$, and the correction is
\begin{equation}
 \frac{B_s^\rho}{\kappa_a}(2p_W-1).
 \label{eq:sg-kappa-correction}
\end{equation}
\end{proposition}
\begin{proof}
The remainder bound forces $R_r=0$ at $U=0$. On the remaining
points, $q_-(y)F_s(y)=p_\rho(y)R_r(y)/B_s^\rho$, with integration
in the continuous case. This proves
equation~\eqref{eq:sg-price-identity}. Good and failure subspaces are
orthogonal. The observable is block diagonal in those subspaces,
so the interference term vanishes and its mean is
$\kappa_a\mathbb E_{q_-}F_s$. The flag probability and decoded
monetary correction follow.
\end{proof}

A failure label in this signed-observable encoding has $W=0$ and
therefore flag probability $1/2$. A different implementation which
accepts only a good preparation and the price flag has probability
$p_{\rm joint}=\kappa_a(1+\mathbb E_{q_-}F_s)/2$. Its correct
decoding is $B_s^\rho(2p_{\rm joint}-\kappa_a)/\kappa_a$.
Confusing these two flags gives a biased price.
Repeated measured postselection has expected $1/\eta$ trials per
accepted state, but does not by itself supply the unitary and adjoint
required by amplitude estimation. I therefore do not multiply that
expected trial cost by an amplitude-estimation query count without
a coherent implementation.

The apparent quotient in $F_s$ can be removed exactly. Taylor's
integral formula gives
\[
 F_s(y)=
 \mathbb E_{Z\sim{\rm Beta}(1,r)}
       \left[\frac{g_y^{(r)}(1+ZU(y))}{H D_r^+}\right].
\]
The integrand is bounded in absolute value by $(1+ZU)^{-r}\leq1$.
At $U=0$ the displayed expression is still defined; that event has
zero $q_-$ weight. A realised derivative oracle and a finite Beta or
quadrature construction are additional costs, with their own biases.

For a fixed $a$-iteration tape, each of its $2a+1$ preparation or
adjoint occurrences contributes at most $\zeta$ operator error.
Telescoping products of unitary factors gives output vector error
at most
$\tau_a=(2a+1)\zeta+\zeta_{\rm refl}$, where
$\zeta_{\rm refl}$ sums the errors of all reflection and phase uses.
The signed observable mean changes by at most $2\tau_a$.
Thus its monetary contribution is at most
$2B_s^\rho\tau_a/\kappa_0$ when the ideal
$\kappa_a\geq\kappa_0>0$.

For clarity, a complete local scalar envelope can also propagate the
normalisers. Suppose $|\widehat M_--M_-^\rho|\leq e_M$,
$|\widehat\kappa-\kappa_a|\leq e_\kappa$ and
$\widehat\kappa\geq\kappa_0>0$. Let $\widehat z\in[-1,1]$ estimate
the signed mean $\zeta_0=\kappa_a\mathbb E_{q_-}F_s$ with radius
$e_z$ at the stated confidence. The decoded correction
$b_s\widehat M_-\widehat z/\widehat\kappa$ has radius at most
\begin{equation}
 \frac{b_s}{\kappa_0}
      \{e_M+M_-^\rho(e_z+e_\kappa)\}.
 \label{eq:sg-scalar-envelope}
\end{equation}
To prove this, split the difference into the moment error, the
signed-mean error and the denominator error. The moment term is
at most $b_se_M/\kappa_0$. The mean term is at most
$b_sM_-^\rho e_z/\kappa_0$. Finally
$|\zeta_0|\leq\kappa_a$ bounds the denominator term by
$b_sM_-^\rho e_\kappa/\kappa_0$. The quantity $e_z$ must include
signed statistical radius, equal to $2e_p$ for probability radius $e_p$,
the $2\tau_a$ state error, value arithmetic and
finite-flag rounding. Polynomial-price, tail, grid and discount errors
are not included in equation~\eqref{eq:sg-scalar-envelope}.

\subsection{A positive-envelope competitor in both measures}

For $|\alpha|=j\leq r$, put
\[
 a_\alpha^\rho=\binom rj\frac{j!}{\alpha!}w^\alpha
                         e^{d_\alpha}Z_\rho(v_\alpha).
\]
All coefficients are positive, and $M_+^\rho=\sum_\alpha
a_\alpha^\rho$. In the continuous case completion of the square gives
\begin{equation}
 q_+(y)=\frac{(T(y)+1)^r\varphi_p(y)}{M_+^{\gamma_p}}
       =\sum_{|\alpha|\leq r}\frac{a_\alpha^{\gamma_p}}
                 {M_+^{\gamma_p}}\varphi_p(y-v_\alpha).
 \label{eq:sg-positive-gaussian-mixture}
\end{equation}
In the finite case the exact component mass is
$p_{\alpha,y}=p_y e^{v_\alpha\cdot y}/Z_\mu(v_\alpha)$, and
$q_+(y)=\sum_\alpha(a_\alpha^\mu/M_+^\mu)p_{\alpha,y}$.
This component's amplitude vector is
$|\varphi_{v_\alpha/2}^\mu\rangle$, not
$|\varphi_{v_\alpha}^\mu\rangle$.
There are $\binom{n+r}{r}$ labelled mixture components before merging.

\begin{proposition}[Positive-mixture rejection]
\label{prop:sg-positive-rejection}
Set $h(y)=((T(y)-1)/(T(y)+1))^r$. Classical rejection from $q_+$
with acceptance probability $h$ has success
$\eta_+=M_-^\rho/M_+^\rho$ and accepted law $q_-$.
A quantum preparation retaining the mixture label,
\[
 |\Omega_+\rangle=\sum_\alpha
       \sqrt{a_\alpha^\rho/M_+^\rho}\,
               |\alpha\rangle|\psi_\alpha^\rho\rangle,
\]
with squared component amplitudes $p_{\alpha,y}$, has financial
marginal $q_+$. A controlled acceptance amplitude $\sqrt h$
gives success $\eta_+$ and financial marginal $q_-$ on success.
\end{proposition}
\begin{proof}
Multiplying $q_+$ by $h$ gives $p_\rho U^r/M_+^\rho$.
Its integral or finite sum is $M_-^\rho/M_+^\rho$; normalising
gives $q_-$. In the quantum purification, distinct component labels
are orthogonal, so their probabilities add to $q_+$ without
interference. The acceptance flag multiplies each amplitude by
$\sqrt h$, producing the same unnormalised financial marginal.
Conditioning or state-specific amplification retains the good vector's
mixture label; a diagonal financial observable has the correct
$q_-$ expectation without erasing that label.
\end{proof}

The expected classical number of proposal trials per accepted sample
is $1/\eta_+$: its survival probabilities are $(1-\eta_+)^j$,
whose geometric sum is $1/\eta_+$. Quantum amplification admits the
$\Theta(\eta_+^{-1/2})$ occurrence order just proved, using the
proposal-and-acceptance unitary with its complete retained label.
The acceptance arithmetic, mixture weights, component preparations,
reflections and their precision must be charged. This is the established
coherent rejection mechanism \citep{ozols2013}; the displayed financial
proposal and its analytic normaliser are the project-specific choice.

For the labelled, unmerged signed expansion,
use $\beta_\alpha^{\gamma_p}=\beta_\alpha$ and continuous component
$|\varphi_\alpha^{\gamma_p}\rangle=\psi_{2v_\alpha}$; the finite
component is $|\varphi_\alpha^\mu\rangle=
|\varphi_{v_\alpha}^\mu\rangle$.
Replacing every sign of $\beta_\alpha^\rho$ by a positive sign gives
$(T+1)^s\sqrt{p_\rho}$. Hence, in either measure,
\begin{equation}
 M_+^\rho=
   \left\|\sum_{\alpha\in\mathcal I_s}
             |\beta_\alpha^\rho|\,|\varphi_\alpha^\rho\rangle\right\|^2
       \leq (L_s^\rho)^2.
 \label{eq:sg-positive-envelope-comparison}
\end{equation}
The inequality is the triangle inequality for unit vectors.
Therefore $\eta_+\geq\eta_{\rm LCU}$, and the positive proposal has
no larger amplification factor in the abstract trial-occurrence
comparison. It need not have fewer gates: it contains degrees through
$2s$, rather than $s$, and computes a bounded ratio for acceptance.
Exact merging of identical signed components can reduce $L_s^\rho$.
Equation~\eqref{eq:sg-positive-envelope-comparison} is not asserted for
that reduced one-norm or for an optimised alternative representation.

The classical opponent should receive the same moments and polynomial
control. For example, $N$ independent accepted $q_-$ samples of
$F_s\in[-1,1]$ obey
$\Pr\{|\overline F_s-\mathbb E F_s|>t\}\leq
2e^{-Nt^2/2}$ by the bounded-variable exponential inequality.
Choosing $t=\varepsilon/B_s^\rho$ gives a sufficient
$N=\lceil2(B_s^\rho/\varepsilon)^2\log(2/\alpha)\rceil$
at failure probability $\alpha$.
The expected proposal count is $NM_+^\rho/M_-^\rho$, with
coefficient preparation and per-trial arithmetic still additional.
This is a concrete equally informed comparator, not a claim that
it is the strongest classical pricer.

\subsection{Weak-dispersion constants with explicit nondegeneracy}

\begin{proposition}[Matched weak-dispersion scaling]
\label{prop:sg-weak-dispersion}
Keep $n,p,w_i,s,K,H,\lambda$ and the base law fixed.
For $0<\delta\leq\delta_0$, assume
\[
 d_i(\delta)=\delta a_i+\delta^2 c_i(\delta),\qquad
 v_i(\delta)=\delta b_i+\delta^2 e_i(\delta),
\]
where all $c_i,e_i$ are uniformly bounded and the centring identities
\eqref{eq:sg-centering} hold for every $\delta$.
Define
$Q_2(y)=\tfrac12\sum_iw_i(a_i+b_i\cdot y)^2$.
For the continuous Gaussian law assume $Q_2$ is not identically
zero. For a fixed finite law require
$\mathbb E_\mu Q_2^{2s}>0$.
Then
\[
 \frac{M_-^\rho}{\delta^{4s}}\longrightarrow
       \mathbb E_\rho Q_2^{2s}>0,\qquad
 L_s^\rho\longrightarrow2^s,\qquad
 M_+^\rho\longrightarrow2^{2s}.
\]
In particular $M_-^\rho=\Theta(\delta^{4s})$ and the two displayed
loading success probabilities are $\Theta(\delta^{4s})$.
On any fixed compact support, $U=O(\delta^2)$ uniformly.
\end{proposition}
\begin{proof}
Put $\ell_i=d_i(\delta)+v_i(\delta)\cdot y$.
Exact centring gives
$U=\sum_iw_i(e^{\ell_i}-1-\ell_i)$.
The integral identity
$e^x-1-x=x^2\int_0^1(1-t)e^{tx}\,dt$ proves
$U/\delta^2\to Q_2$ pointwise, and uniformly on compact sets.
There is a constant $C$ such that
$|\ell_i|\leq C\delta(1+\|y\|)$ for every permitted $\delta$.
Thus
\[
 0\leq U/\delta^2\leq
       \tfrac12 C^2(1+\|y\|)^2
                      e^{C\delta_0(1+\|y\|)}.
\]
Every finite power of this bound is Gaussian integrable: after
raising to that power, the exponential linear in $\|y\|$ is
bounded using $a\|y\|\leq\|y\|^2/4+a^2$, leaving a polynomial
times $e^{-\|y\|^2/4}$ in the density.
Dominated convergence therefore proves the moment limit for
$\gamma_p$. A nonzero $Q_2\geq0$ is positive on an open ball;
the Gaussian density is positive there, so the limit is positive.
For a finite law, uniform convergence at finitely many nodes gives
the same limit, and the assumed positive finite moment supplies
nondegeneracy.

Each $d_\alpha,v_\alpha$ tends to zero. In the continuous formula
$e^{d_\alpha+\|v_\alpha\|^2}\to1$; in the finite formula
$e^{d_\alpha}\sqrt{Z_\mu(2v_\alpha)}\to1$.
Consequently
\[
 L_s^\rho\to\sum_{j=0}^s\binom sj
      \sum_{|\alpha|=j}\frac{j!}{\alpha!}w^\alpha
    =\sum_{j=0}^s\binom sj(\sum_iw_i)^j=2^s.
\]
The same Gaussian exponential envelope or a finite sum gives
$M_+^\rho\to\mathbb E_\rho2^{2s}=2^{2s}$.
The positive limits prove all the stated Theta conclusions.
\end{proof}

A sufficient GBM subfamily has equal initial stock levels and equal
dividend yields, a fixed common loading $\lambda$, and residual
loadings $\delta u_i$. Its risk-neutral deterministic contrasts are
of order $\delta^2$, from the volatility compensators, and its
centred residual contrasts are
$\delta(u_i-\sum_jw_j u_j)$. Thus $a_i=0$ and the proposition
applies when those contrasts are not all zero.
Fixed unequal initial levels or dividend yields can leave a nonzero
contrast as $\delta\to0$ and do not automatically satisfy this
weak-dispersion family. Scaling the common smoothing $\lambda$
with $\delta$ also changes $b_s$ and invalidates a comparison that
treats it as fixed.

For the same fixed finite law $\mu$, choose a certified support upper
bound $U_*(\delta)=C_*\delta^2$ with $C_*>0$.
The preceding compact bound supplies such a constant.
Nondegeneracy on a positive-mass grid point also makes the exact
maximum $\Theta(\delta^2)$.
The direct signed remainder can be normalised by
$\rho_s=b_sU_*^{2s}$, so it has range contained in $[-1,1]$.
The integral derivative representation supplies this observable
without a division by $U$.

Let $\varepsilon$ denote the monetary allowance remaining for the
correction after the separately charged deterministic errors.
Consider the active monetary regime
$\varepsilon=o(\delta^{4s})$ at fixed confidence.
The usual worst-case additive-accuracy amplitude-estimation allocation
uses $O(B_s^\mu/\varepsilon)$ invocations of the amplified
importance-and-observable circuit when $\kappa_a$ is bounded below
by a fixed positive constant.
Multiplying by its displayed loading occurrence order gives
\[
 O\!\left(\frac{b_s L_s^\mu\sqrt{M_-^\mu}}{\varepsilon}\right)
       =O(\delta^{2s}/\varepsilon)
\]
for signed LCU, and
$O(b_s\sqrt{M_-^\mu M_+^\mu}/\varepsilon)
=O(\delta^{2s}/\varepsilon)$ for positive-mixture rejection.
The direct range allocation is $O(\rho_s/\varepsilon)
=O(\delta^{4s}/\varepsilon)$.
For the corresponding canonical allocations these displayed budgets
have the stated orders up to fixed constants and confidence factors.
They are not instance-specific lower bounds. An analytic mean,
a stronger variance-sensitive algorithm or a different representation
can require less work.

Both competitors here estimate the same finite correction.
Comparing a full continuous importance law with a finite-support
direct range without the common transfer certificate would change
the estimand. If the financial support grows with tolerance, or
the node precision or derivative order changes with $\delta$, its
range and implementation cost must be analysed again.
The preceding comparison shows a disadvantage of the displayed
importance loaders in this fixed-support regime when evaluator
costs are comparable; it does not exclude an advantage in another
proved regime.

If $B_s^\rho\leq\varepsilon$, the deterministic remainder interval
$[-B_s^\rho,B_s^\rho]$ already resolves the correction and its
midpoint requires no quantum calls. An implementation must check
this before claiming an estimation saving.

\subsection{Accounting and the limits of the retained construction}

The continuous exact identity preserves the original risk-neutral
expectation, followed by the original deterministic discount.
The finite identity preserves its stated finite numerical target,
which requires the separately validated transfer to that model price.
The importance law alone is not a replacement pricing measure.
A no-arbitrage conclusion follows from that original classical
model and the exact identity, not from interpreting computational
amplitudes as new asset dynamics. Independently rounded numerical
prices do not automatically satisfy joint across-contract arbitrage
constraints.

For a complete implementation I must retain the finite base law,
component normalisers, coefficient signs and their arithmetic
enclosures; every preparation and controlled component tape;
relative phases; clean private work; selector and success predicates;
the full zero-state and good-subspace reflections; each adjoint;
amplification choices and a certified $\kappa_a$; the derivative or
remainder evaluator; finite auxiliary laws; the signed monetary
flag; the mean estimator and its confidence postprocessing;
observations and resets. Classical table construction, Gaussian
moment evaluation and coefficient precision are setup costs.
They cannot be assigned zero cost or amortised without a stated
reuse workload. A measured lookup or conditional postselection tape
does not acquire a coherent adjoint by being reversed.

The column and operator error bounds above show why severe
cancellation can increase required component precision. Tight
derivative constants can likewise introduce their own shared
generation cost. The bounded integral observable removes the
small-$U$ quotient but does not prove that its reversible
Gaussian-density, logarithm and polynomial arithmetic is cheap.
The retained finite toy tests only a preparation signal and its
amplification. It does not implement all of these financial and
statistical operations. I therefore retain these identities and
comparators as a supplementary specification and a rejected-loading
analysis, without a complete pricing-cost or quantum-advantage claim.
